\documentclass[border=4pt]{standalone}
\usepackage{amsmath,amssymb,mathtools,xcolor,varwidth}
\definecolor{radgreen}{HTML}{00A35A}
\definecolor{circblue}{HTML}{1E6FE0}
\definecolor{centerorange}{HTML}{E8770A}
\begin{document}
\begin{varwidth}{28em}
\large\bfseries
Proposition VII is the first concrete application of the force-measure proved in \textcolor{radgreen}{Proposition VI},      now applied to the simplest curvilinear orbit: the \textcolor{circblue}{circle $VQPA$}.      A body $P$ revolves in the circumference and $Q$ is the next place into which it moves.      The remarkable freedom here is that the \textcolor{centerorange}{force centre $S$} need not sit at the centre $O$ —      it may lie anywhere inside the circle, and we still seek the law along the \textcolor{radgreen}{radius $SP$},      which is the distance whose powers Prop. VI taught us to track.
\end{varwidth}
\end{document}
